% ---------------------------------------------------------------------------
% §3.2 Empirical Strategy — "The second stage is a survey experiment." ile
% başlayan paragrafın YERİNE geçecek metin ("burayı daha sonra yazalim"
% notunun olduğu yer).
% ---------------------------------------------------------------------------

The second stage is a survey experiment conducted on the synthetic
population. Each calibrated digital twin is exposed to two versions of the
same protest event, embedded in short news-style vignettes that vary only
the political structure surrounding the event: a \emph{threat} condition, in
which the march for Kurdish-language education takes place against the
backdrop of an active cross-border military operation, and an
\emph{opportunity} condition, in which the same march takes place during an
ongoing negotiation process. Because the manipulation operates on the
political structure of the moment rather than on the content of the claim,
the design speaks directly to the political-opportunity and threat framework
developed in Section~2. Every twin answers in both conditions in independent
prompt turns, which yields a within-persona counterfactual contrast that no
conventional survey design can produce: the same respondent observed under
two political structures. Treatment effects are evaluated as aggregate
distributional shifts between conditions, consistent with the calibration
target of the first stage. The full design is described in Section~5.
